\chapter{Conclusioni}
\label{cap:conclusioni}

\section{Prodotto allo stato attuale}
Il prodotto allo stato attuale permette il testing dei modelli di intelligenza artificiale generativa rispetto alle vulnerabilità OWASP top 10 per AI generativa.
\subsection{Dashboard}
L'utilizzo del prodotto è semplice e intuitivo, grazie alla dashboard web che consente di visualizzare i risultati dei test in modo chiaro e dettagliato.
Al primo accesso l'utente viene accolto dalla dashboard che presenta una panoramica dei test eseguiti, con informazioni chiave come nome del test, target, dataset utilizzato, status, coverage e età del test.
Da qui l'utente può avviare un test o accedere ai dettagli di ogni test, visualizzando informazioni più approfondite e i risultati specifici per ogni categoria di vulnerabilità OWASP. La dashboard offre anche grafici e visualizzazioni che facilitano l'interpretazione dei risultati, permettendo all'utente di identificare rapidamente le aree di miglioramento per i modelli testati.


\begin{figure}[!h] 
    \centering 
    \includegraphics[width=0.9\textwidth]{dashboard.png} 
    \caption{Dashboard del prodotto}
\end{figure}
\FloatBarrier
\clearpage
\subsection{Avvio del test}
L'avvio di un test è semplice e guidato. L'utente seleziona il target di test, il dataset da utilizzare e lo scorer per valutare i risultati. Una volta avviato il test, il sistema esegue automaticamente una serie di attacchi simulati contro il modello di intelligenza artificiale generativa, analizzando le risposte del modello e identificando eventuali vulnerabilità. I risultati vengono poi visualizzati nella dashboard, consentendo all'utente di esaminare i dettagli di ogni test e di prendere decisioni informate per migliorare la sicurezza dei modelli testati. L'interfaccia è progettata per essere user-friendly, con istruzioni chiare e opzioni configurabili per adattarsi alle esigenze specifiche dell'utente e del progetto.

\begin{figure}[!h] 
    \centering 
    \includegraphics[width=0.7\textwidth]{newRun.png} 
    \caption{Avvio del test}
\end{figure}
\FloatBarrier


\clearpage
\subsection{Dettagli del test}

Nella sezione dettagli del test, l'utente può visualizzare informazioni specifiche sui risultati dei test, inclusi i prompt utilizzati, i payload degli attacchi e le risposte del modello. Questa sezione fornisce una visione approfondita delle vulnerabilità identificate, consentendo all'utente di comprendere meglio le debolezze del modello e di pianificare interventi correttivi mirati. La visualizzazione dei dettagli del test è fondamentale per l'analisi approfondita dei risultati e per la definizione di strategie di mitigazione efficaci.

\begin{figure}[!h] 
    \centering 
    \includegraphics[width=0.9\textwidth]{details.png} 
    \caption{Dettagli del test}
\end{figure}
\FloatBarrier

\begin{figure}[!h] 
    \centering 
    \begin{minipage}{0.48\textwidth}
        \centering
        \includegraphics[width=\textwidth]{esempioPassed.png}
        \caption{Esempio test superato}
    \end{minipage}
    \hfill
    \begin{minipage}{0.48\textwidth}
        \centering
        \includegraphics[width=\textwidth]{esempioFailed.png}
        \caption{Esempio test fallito}
    \end{minipage}
\end{figure}
\FloatBarrier
\clearpage

\subsection{Pagina di documentazione}
La pagina di documentazione fornisce all'utente tutte le informazioni necessarie per comprendere il funzionamento del sistema, le tecnologie utilizzate e le istruzioni per l'uso. Questa sezione è essenziale per garantire che gli utenti possano sfruttare appieno le funzionalità del prodotto e per facilitare la risoluzione di eventuali problemi o dubbi che potrebbero sorgere durante l'utilizzo. La documentazione è strutturata in modo chiaro e accessibile, con sezioni dedicate a diversi aspetti del sistema, come l'architettura, le funzionalità principali e le best practice per l'utilizzo efficace del prodotto.

\begin{figure}[!h] 
    \centering 
    \includegraphics[width=0.9\textwidth]{documentation.png} 
    \caption{Pagina di documentazione}
\end{figure}
\FloatBarrier
\clearpage

\section{Stato dei requisiti}
Di seguito viene riportata una tabella riassuntiva dello stato di soddisfacimento dei requisiti al termine del progetto di stage.

\begin{table}[htbp]
\centering
\caption{Stato di soddisfacimento dei requisiti al termine dello stage}
\label{tab:stato-requisiti}
\begin{tabularx}{\textwidth}{lXc}
\hline\hline
\textbf{Requisito} & \textbf{Descrizione} & \textbf{Soddisfatto}\\
\hline
RFN-1 & Visualizzazione dashboard principale con statistiche dei test & Sì \\
RFN-2 & Visualizzazione card con statistiche aggregate dei test & Sì \\
RFN-3 & Visualizzazione grafico temporale andamento test & Sì \\
RFN-4 & Visualizzazione tabella test eseguiti sui target & Sì \\
RFN-5 & Visualizzazione dettagli singolo test (nome, scorer, dataset, target, status, coverage, età) & Sì \\
\hline
RQD-1 & Architettura modulare per estensione funzionalità (nuovi dataset, nuovi scorer) & Sì \\
RQD-2 & Interfaccia dashboard intuitiva senza competenze tecniche avanzate & Sì \\
RQD-3 & Gestione efficiente dataset di grandi dimensioni con feedback avanzamento & Sì \\
RQD-4 & Codice commentato e documentato per manutenzione e handover & Sì \\
RQD-5 & Gestione errori imprevisti API con messaggi chiari & Sì \\
RQD-6 & Scalabilità per test paralleli e dataset grandi & Parziale \\
RQD-7 & Interfaccia responsive per desktop e laptop & Sì \\
\hline
RVO-1 & Compatibilità con infrastrutture cloud/locali Var Group & Sì \\
RVO-2 & Controlli per vulnerabilità OWASP Top 10 for LLM & Sì \\
RVO-3 & Accessibilità tramite browser moderni senza plugin & Sì \\
RVO-4 & Integrazione modelli AI tramite API standard & Sì \\
\hline\hline
\end{tabularx}
\end{table}

Come evidenziato dalla tabella, la maggior parte dei requisiti è stata soddisfatta. L'unico requisito parzialmente soddisfatto riguarda la scalabilità per l'esecuzione di test in parallelo (RQD-6), funzionalità che richiederebbe ulteriore sviluppo per essere completamente implementata.
\clearpage

\section{Valutazione personale}
In conclusione il progetto di Stage mi ha permesso di acquisire competenze specifiche su strumenti e tecnologie all'avanguardia utilizzati nel mondo del lavoro.
Nonostante le difficoltà tecniche riscontrate durante lo sviluppo, il progetto è stato portato a termine con successo, raggiungendo gli obiettivi prefissati. Un possibile miglioramento futuro riguarderebbe la semplificazione del cambio del target di test, per renderlo più immediato e flessibile per l'utente finale.

Il lavoro di sviluppo del prodotto mi ha dato l'opportunità di applicare concretamente le conoscenze acquisite durante il percorso universitario, affrontando sfide reali e imparando a risolverle in modo efficace. Ho potuto approfondire la mia comprensione dei modelli di linguaggio e delle tecniche di testing, oltre a migliorare le mie capacità di programmazione e problem-solving.
L'esperienza di lavorare in un ambiente professionale mi ha insegnato l'importanza della collaborazione, della comunicazione e della gestione del tempo, competenze fondamentali per il successo in qualsiasi carriera. Sono grato per l'opportunità di crescita personale e professionale che Var Group mi ha offerto durante questo stage, e sono fiducioso che le competenze acquisite mi saranno preziose per il mio futuro professionale.
